\documentclass[../DoAn.tex]{subfiles}
\begin{document}

\begin{center}
    \Large{\textbf{TÓM TẮT NỘI DUNG ĐỒ ÁN}}\\
\end{center}
\vspace{1cm}
Video bài giảng là hình thức học tập trực tuyến phổ biến nhờ chi phí thấp và khả năng tiếp cận rộng rãi. Tuy nhiên, do mang bản chất một chiều, hình thức này dễ khiến người học rơi vào trạng thái thụ động và đánh giá sai mức độ hiểu bài, trong khi việc kiểm tra mức độ tiếp thu thường bị tách rời khỏi quá trình xem video. Các nền tảng hiện nay hoặc chưa hỗ trợ hiệu quả các hoạt động tương tác ngay trong bài giảng, hoặc vẫn yêu cầu giảng viên tự thiết kế thủ công từng hoạt động, gây tốn nhiều công sức.

Trước hạn chế đó, đồ án lựa chọn hướng khai thác trí tuệ nhân tạo, đặc biệt là các mô hình ngôn ngữ lớn có khả năng hiểu nội dung và tự đề xuất câu hỏi bám sát bài giảng, để tự động sinh nội dung tương tác từ video. Cách tiếp cận này giúp giảm gánh nặng soạn thảo cho giảng viên đồng thời tăng tính chủ động cho người học.

Trên cơ sở đó, đồ án xây dựng Dyadia, một nền tảng học tập qua video tích hợp trí tuệ nhân tạo. Nền tảng cho phép giảng viên tổ chức nội dung theo khóa học và bài giảng, tải lên video, tự động sinh các hoạt động tương tác, và quản lý người học theo vai trò. Về phía người học, mỗi bài giảng trở thành một phiên học có tương tác với chấm điểm và phản hồi tự động ngay trong khi xem; giảng viên đồng thời theo dõi được kết quả học tập của lớp.

Đóng góp chính của đồ án là đề xuất và hiện thực hóa quy trình chuyển hóa video bài giảng thụ động thành bài học chủ động với chi phí soạn thảo thấp. Hệ thống đã được xây dựng đầy đủ luồng nghiệp vụ, từ tải video lên, phân tích nội dung và sinh hoạt động tương tác, tới học tập, chấm điểm, phản hồi tự động và theo dõi kết quả.
\vfill
\begin{flushright}
Sinh viên thực hiện\\
\begin{tabular}{@{}c@{}}
\textit{(Ký và ghi rõ họ tên)}
\end{tabular}
\end{flushright}

\end{document}
